\documentclass[journal,12pt,twocolumn]{IEEEtran}
%
\usepackage{setspace}
\usepackage{gensymb}
\usepackage{xcolor}
\usepackage{caption}
\usepackage[hyphens,spaces,obeyspaces]{url}
%\usepackage{subcaption}
%\doublespacing
\singlespacing

%\usepackage{graphicx}
%\usepackage{amssymb}
%\usepackage{relsize}
\usepackage[cmex10]{amsmath}
\usepackage{mathtools}
%\usepackage{amsthm}
%\interdisplaylinepenalty=2500
%\savesymbol{iint}
%\usepackage{txfonts}
%\restoresymbol{TXF}{iint}
%\usepackage{wasysym}
\usepackage{amsthm}
\usepackage{mathrsfs}
\usepackage{txfonts}
\usepackage{stfloats}
\usepackage{cite}
\usepackage{cases}
\usepackage{subfig}
%\usepackage{xtab}
\usepackage{longtable}
\usepackage{multirow}
%\usepackage{algorithm}
%\usepackage{algpseudocode}
\usepackage{enumerate}
\usepackage{mathtools}
%\usepackage{eenrc}
%\usepackage[framemethod=tikz]{mdframed}
\usepackage[breaklinks]{hyperref}
%\usepackage{breakcites}
\usepackage{listings}
\usepackage[latin1]{inputenc}                                 %%
\usepackage{color}                                            %%
\usepackage{array}                                            %%
\usepackage{longtable}                                        %%
\usepackage{calc}                                             %%
\usepackage{multirow}                                         %%
\usepackage{hhline}                                           %%
\usepackage{ifthen}                                           %%
%optionally (for landscape tables embedded in another document): %%
\usepackage{lscape}     
\usepackage{amsmath}
\usepackage{tikz}
\usepackage{circuitikz}
\usepackage{karnaugh-map}
\usepackage{pgf}
\usepackage[hyphenbreaks]{breakurl}
\usepackage{enumitem}

%\usepackage{url}
%\def\UrlBreaks{\do\/\do-}





%\usepackage{stmaryrd}


%\usepackage{wasysym}
%\newcounter{MYtempeqncnt}
\DeclareMathOperator*{\Res}{Res}
%\renewcommand{\baselinestretch}{2}
\renewcommand\thesection{\arabic{section}}
\renewcommand\thesubsection{\thesection.\arabic{subsection}}
\renewcommand\thesubsubsection{\thesubsection.\arabic{subsubsection}}

\renewcommand\thesectiondis{\arabic{section}}
\renewcommand\thesubsectiondis{\thesectiondis.\arabic{subsection}}
\renewcommand\thesubsubsectiondis{\thesubsectiondis.\arabic{subsubsection}}

% correct bad hyphenation here
\hyphenation{op-tical net-works semi-conduc-tor}

%\lstset{
	%language=C,
	%frame=single, 
	%breaklines=true
	%}

%\lstset{
	%%basicstyle=\small\ttfamily\bfseries,
	%%numberstyle=\small\ttfamily,
	%language=Octave,
	%backgroundcolor=\color{white},
	%%frame=single,
	%%keywordstyle=\bfseries,
	%%breaklines=true,
	%%showstringspaces=false,
	%%xleftmargin=-10mm,
	%%aboveskip=-1mm,
	%%belowskip=0mm
	%}

%\surroundwithmdframed[width=\columnwidth]{lstlisting}
\def\inputGnumericTable{}                                 %%
\lstset{
	%language=C,
	frame=single, 
	breaklines=true,
	columns=fullflexible
}

\graphicspath{ {./title/} } 
\begin{document}
	
	
	\title{
		
		
		\includegraphics[width=15cm, height=4cm]{title}
		\centering
		
	}
	\author{G V V Sharma$^{*}$% <-this % stops a space
		\thanks{*The author is with the Department
			of Electrical Engineering, Indian Institute of Technology, Hyderabad
			502285 India e-mail:  gadepall@iith.ac.in. All content in this manual is released under GNU GPL.  Free and open source.}}

% make the title area
\maketitle

\tableofcontents

\bigskip

\renewcommand{\thefigure}{\theenumi}
\renewcommand{\thetable}{\theenumi}

\begin{abstract}
%\boldmath
This manual explains Karnaugh maps (K-map) by finding the
logic functions for the incrementing decoder.
\end{abstract}

%
\section{Incrementing Decoder}
The incrementing decoder   takes the numbers $0,1,\dots,9$ in binary as inputs and generates
the consecutive number as output.  The corresponding truth table is available in Table. \ref{table:counter_decoder}.
\begin{table}[!h]
	\centering	
	%%%%%%%%%%%%%%%%%%%%%%%%%%%%%%%%%%%%%%%%%%%%%%%%%%%%%%%%%%%%%%%%%%%%%%
%%                                                                  %%
%%  This is the header of a LaTeX2e file exported from Gnumeric.    %%
%%                                                                  %%
%%  This file can be compiled as it stands or included in another   %%
%%  LaTeX document. The table is based on the longtable package so  %%
%%  the longtable options (headers, footers...) can be set in the   %%
%%  preamble section below (see PRAMBLE).                           %%
%%                                                                  %%
%%  To include the file in another, the following two lines must be %%
%%  in the including file:                                          %%
%%        \def\inputGnumericTable{}                                 %%
%%  at the beginning of the file and:                               %%
%%        \input{name-of-this-file.tex}                             %%
%%  where the table is to be placed. Note also that the including   %%
%%  file must use the following packages for the table to be        %%
%%  rendered correctly:                                             %%
%%    \usepackage[latin1]{inputenc}                                 %%
%%    \usepackage{color}                                            %%
%%    \usepackage{array}                                            %%
%%    \usepackage{longtable}                                        %%
%%    \usepackage{calc}                                             %%
%%    \usepackage{multirow}                                         %%
%%    \usepackage{hhline}                                           %%
%%    \usepackage{ifthen}                                           %%
%%  optionally (for landscape tables embedded in another document): %%
%%    \usepackage{lscape}                                           %%
%%                                                                  %%
%%%%%%%%%%%%%%%%%%%%%%%%%%%%%%%%%%%%%%%%%%%%%%%%%%%%%%%%%%%%%%%%%%%%%%



%%  This section checks if we are begin input into another file or  %%
%%  the file will be compiled alone. First use a macro taken from   %%
%%  the TeXbook ex 7.7 (suggestion of Han-Wen Nienhuys).            %%
\def\ifundefined#1{\expandafter\ifx\csname#1\endcsname\relax}


%%  Check for the \def token for inputed files. If it is not        %%
%%  defined, the file will be processed as a standalone and the     %%
%%  preamble will be used.                                          %%
\ifundefined{inputGnumericTable}

%%  We must be able to close or not the document at the end.        %%
	\def\gnumericTableEnd{\end{document}}


%%%%%%%%%%%%%%%%%%%%%%%%%%%%%%%%%%%%%%%%%%%%%%%%%%%%%%%%%%%%%%%%%%%%%%
%%                                                                  %%
%%  This is the PREAMBLE. Change these values to get the right      %%
%%  paper size and other niceties.                                  %%
%%                                                                  %%
%%%%%%%%%%%%%%%%%%%%%%%%%%%%%%%%%%%%%%%%%%%%%%%%%%%%%%%%%%%%%%%%%%%%%%

	\documentclass[12pt%
			  %,landscape%
                    ]{report}
       \usepackage[latin1]{inputenc}
       \usepackage{fullpage}
       \usepackage{color}
       \usepackage{array}
       \usepackage{longtable}
       \usepackage{calc}
       \usepackage{multirow}
       \usepackage{hhline}
       \usepackage{ifthen}

	\begin{document}


%%  End of the preamble for the standalone. The next section is for %%
%%  documents which are included into other LaTeX2e files.          %%
\else

%%  We are not a stand alone document. For a regular table, we will %%
%%  have no preamble and only define the closing to mean nothing.   %%
    \def\gnumericTableEnd{}

%%  If we want landscape mode in an embedded document, comment out  %%
%%  the line above and uncomment the two below. The table will      %%
%%  begin on a new page and run in landscape mode.                  %%
%       \def\gnumericTableEnd{\end{landscape}}
%       \begin{landscape}


%%  End of the else clause for this file being \input.              %%
\fi

%%%%%%%%%%%%%%%%%%%%%%%%%%%%%%%%%%%%%%%%%%%%%%%%%%%%%%%%%%%%%%%%%%%%%%
%%                                                                  %%
%%  The rest is the gnumeric table, except for the closing          %%
%%  statement. Changes below will alter the table's appearance.     %%
%%                                                                  %%
%%%%%%%%%%%%%%%%%%%%%%%%%%%%%%%%%%%%%%%%%%%%%%%%%%%%%%%%%%%%%%%%%%%%%%

\providecommand{\gnumericmathit}[1]{#1} 
%%  Uncomment the next line if you would like your numbers to be in %%
%%  italics if they are italizised in the gnumeric table.           %%
%\renewcommand{\gnumericmathit}[1]{\mathit{#1}}
\providecommand{\gnumericPB}[1]%
{\let\gnumericTemp=\\#1\let\\=\gnumericTemp\hspace{0pt}}
 \ifundefined{gnumericTableWidthDefined}
        \newlength{\gnumericTableWidth}
        \newlength{\gnumericTableWidthComplete}
        \newlength{\gnumericMultiRowLength}
        \global\def\gnumericTableWidthDefined{}
 \fi
%% The following setting protects this code from babel shorthands.  %%
 \ifthenelse{\isundefined{\languageshorthands}}{}{\languageshorthands{english}}
%%  The default table format retains the relative column widths of  %%
%%  gnumeric. They can easily be changed to c, r or l. In that case %%
%%  you may want to comment out the next line and uncomment the one %%
%%  thereafter                                                      %%
\providecommand\gnumbox{\makebox[0pt]}
%%\providecommand\gnumbox[1][]{\makebox}

%% to adjust positions in multirow situations                       %%
\setlength{\bigstrutjot}{\jot}
\setlength{\extrarowheight}{\doublerulesep}

%%  The \setlongtables command keeps column widths the same across  %%
%%  pages. Simply comment out next line for varying column widths.  %%
\setlongtables

\setlength\gnumericTableWidth{%
	13pt+%
	12pt+%
	12pt+%
	12pt+%
	12pt+%
	12pt+%
	12pt+%
	12pt+%
0pt}
\def\gumericNumCols{8}
\setlength\gnumericTableWidthComplete{\gnumericTableWidth+%
         \tabcolsep*\gumericNumCols*2+\arrayrulewidth*\gumericNumCols}
\ifthenelse{\lengthtest{\gnumericTableWidthComplete > \linewidth}}%
         {\def\gnumericScale{\ratio{\linewidth-%
                        \tabcolsep*\gumericNumCols*2-%
                        \arrayrulewidth*\gumericNumCols}%
{\gnumericTableWidth}}}%
{\def\gnumericScale{1}}

%%%%%%%%%%%%%%%%%%%%%%%%%%%%%%%%%%%%%%%%%%%%%%%%%%%%%%%%%%%%%%%%%%%%%%
%%                                                                  %%
%% The following are the widths of the various columns. We are      %%
%% defining them here because then they are easier to change.       %%
%% Depending on the cell formats we may use them more than once.    %%
%%                                                                  %%
%%%%%%%%%%%%%%%%%%%%%%%%%%%%%%%%%%%%%%%%%%%%%%%%%%%%%%%%%%%%%%%%%%%%%%

\ifthenelse{\isundefined{\gnumericColA}}{\newlength{\gnumericColA}}{}\settowidth{\gnumericColA}{\begin{tabular}{@{}p{13pt*\gnumericScale}@{}}x\end{tabular}}
\ifthenelse{\isundefined{\gnumericColB}}{\newlength{\gnumericColB}}{}\settowidth{\gnumericColB}{\begin{tabular}{@{}p{13pt*\gnumericScale}@{}}x\end{tabular}}
\ifthenelse{\isundefined{\gnumericColC}}{\newlength{\gnumericColC}}{}\settowidth{\gnumericColC}{\begin{tabular}{@{}p{13pt*\gnumericScale}@{}}x\end{tabular}}
\ifthenelse{\isundefined{\gnumericColD}}{\newlength{\gnumericColD}}{}\settowidth{\gnumericColD}{\begin{tabular}{@{}p{13pt*\gnumericScale}@{}}x\end{tabular}}
\ifthenelse{\isundefined{\gnumericColE}}{\newlength{\gnumericColE}}{}\settowidth{\gnumericColE}{\begin{tabular}{@{}p{13pt*\gnumericScale}@{}}x\end{tabular}}
\ifthenelse{\isundefined{\gnumericColF}}{\newlength{\gnumericColF}}{}\settowidth{\gnumericColF}{\begin{tabular}{@{}p{13pt*\gnumericScale}@{}}x\end{tabular}}
\ifthenelse{\isundefined{\gnumericColG}}{\newlength{\gnumericColG}}{}\settowidth{\gnumericColG}{\begin{tabular}{@{}p{13pt*\gnumericScale}@{}}x\end{tabular}}
\ifthenelse{\isundefined{\gnumericColH}}{\newlength{\gnumericColH}}{}\settowidth{\gnumericColH}{\begin{tabular}{@{}p{13pt*\gnumericScale}@{}}x\end{tabular}}


\begin{tabular}[c]{
%\begin{longtable}[c]{%
	b{\gnumericColA}%
	b{\gnumericColB}%
	b{\gnumericColC}%
	b{\gnumericColD}%
	b{\gnumericColE}%
	b{\gnumericColF}%
	b{\gnumericColG}%
	b{\gnumericColH}%
	}

%%%%%%%%%%%%%%%%%%%%%%%%%%%%%%%%%%%%%%%%%%%%%%%%%%%%%%%%%%%%%%%%%%%%%%
%%  The longtable options. (Caption, headers... see Goosens, p.124) %%
%	\caption{The Table Caption.}             \\	%
% \hline	% Across the top of the table.
%%  The rest of these options are table rows which are placed on    %%
%%  the first, last or every page. Use \multicolumn if you want.    %%

%%  Header for the first page.                                      %%
%	\multicolumn{8}{c}{The First Header} \\ \hline 
%	\multicolumn{1}{c}{colTag}	%Column 1
%	&\multicolumn{1}{c}{colTag}	%Column 2
%	&\multicolumn{1}{c}{colTag}	%Column 3
%	&\multicolumn{1}{c}{colTag}	%Column 4
%	&\multicolumn{1}{c}{colTag}	%Column 5
%	&\multicolumn{1}{c}{colTag}	%Column 6
%	&\multicolumn{1}{c}{colTag}	%Column 7
%	&\multicolumn{1}{c}{colTag}	\\ \hline %Last column
%	\endfirsthead

%%  The running header definition.                                  %%
%	\hline
%	\multicolumn{8}{l}{\ldots\small\slshape continued} \\ \hline
%	\multicolumn{1}{c}{colTag}	%Column 1
%	&\multicolumn{1}{c}{colTag}	%Column 2
%	&\multicolumn{1}{c}{colTag}	%Column 3
%	&\multicolumn{1}{c}{colTag}	%Column 4
%	&\multicolumn{1}{c}{colTag}	%Column 5
%	&\multicolumn{1}{c}{colTag}	%Column 6
%	&\multicolumn{1}{c}{colTag}	%Column 7
%	&\multicolumn{1}{c}{colTag}	\\ \hline %Last column
%	\endhead

%%  The running footer definition.                                  %%
%	\hline
%	\multicolumn{8}{r}{\small\slshape continued\ldots} \\
%	\endfoot

%%  The ending footer definition.                                   %%
%	\multicolumn{8}{c}{That's all folks} \\ \hline 
%	\endlastfoot
%%%%%%%%%%%%%%%%%%%%%%%%%%%%%%%%%%%%%%%%%%%%%%%%%%%%%%%%%%%%%%%%%%%%%%

\hhline{|-|-|-|-|-|-|-|-}
	 \multicolumn{1}{|p{\gnumericColA}|}%
	{\gnumericPB{\centering}\gnumbox{Z}}
	&\multicolumn{1}{p{\gnumericColB}|}%
	{\gnumericPB{\centering}\gnumbox{Y}}
	&\multicolumn{1}{p{\gnumericColC}|}%
	{\gnumericPB{\centering}\gnumbox{X}}
	&\multicolumn{1}{p{\gnumericColD}|}%
	{\gnumericPB{\centering}\gnumbox{W}}
	&\multicolumn{1}{p{\gnumericColE}|}%
	{\gnumericPB{\centering}\gnumbox{\textbf{D}}}
	&\multicolumn{1}{p{\gnumericColF}|}%
	{\gnumericPB{\centering}\gnumbox{\textbf{C}}}
	&\multicolumn{1}{p{\gnumericColG}|}%
	{\gnumericPB{\centering}\gnumbox{\textbf{B}}}
	&\multicolumn{1}{p{\gnumericColH}|}%
	{\gnumericPB{\centering}\gnumbox{\textbf{A}}}
\\
\hhline{|--------|}
	 \multicolumn{1}{|p{\gnumericColA}|}%
	{\gnumericPB{\centering}\gnumbox{0}}
	&\multicolumn{1}{p{\gnumericColB}|}%
	{\gnumericPB{\centering}\gnumbox{0}}
	&\multicolumn{1}{p{\gnumericColC}|}%
	{\gnumericPB{\centering}\gnumbox{0}}
	&\multicolumn{1}{p{\gnumericColD}|}%
	{\gnumericPB{\centering}\gnumbox{0}}
	&\multicolumn{1}{p{\gnumericColE}|}%
	{\gnumericPB{\centering}\gnumbox{\textbf{1}}}
	&\multicolumn{1}{p{\gnumericColF}|}%
	{\gnumericPB{\centering}\gnumbox{\textbf{0}}}
	&\multicolumn{1}{p{\gnumericColG}|}%
	{\gnumericPB{\centering}\gnumbox{\textbf{0}}}
	&\multicolumn{1}{p{\gnumericColH}|}%
	{\gnumericPB{\centering}\gnumbox{\textbf{1}}}
\\
\hhline{|--------|}
	 \multicolumn{1}{|p{\gnumericColA}|}%
	{\gnumericPB{\centering}\gnumbox{1}}
	&\multicolumn{1}{p{\gnumericColB}|}%
	{\gnumericPB{\centering}\gnumbox{0}}
	&\multicolumn{1}{p{\gnumericColC}|}%
	{\gnumericPB{\centering}\gnumbox{0}}
	&\multicolumn{1}{p{\gnumericColD}|}%
	{\gnumericPB{\centering}\gnumbox{1}}
	&\multicolumn{1}{p{\gnumericColE}|}%
	{\gnumericPB{\centering}\gnumbox{\textbf{1}}}
	&\multicolumn{1}{p{\gnumericColF}|}%
	{\gnumericPB{\centering}\gnumbox{\textbf{0}}}
	&\multicolumn{1}{p{\gnumericColG}|}%
	{\gnumericPB{\centering}\gnumbox{\textbf{0}}}
	&\multicolumn{1}{p{\gnumericColH}|}%
	{\gnumericPB{\centering}\gnumbox{\textbf{0}}}
\\
\hhline{|--------|}
	 \multicolumn{1}{|p{\gnumericColA}|}%
	{\gnumericPB{\centering}\gnumbox{1}}
	&\multicolumn{1}{p{\gnumericColB}|}%
	{\gnumericPB{\centering}\gnumbox{0}}
	&\multicolumn{1}{p{\gnumericColC}|}%
	{\gnumericPB{\centering}\gnumbox{0}}
	&\multicolumn{1}{p{\gnumericColD}|}%
	{\gnumericPB{\centering}\gnumbox{0}}
	&\multicolumn{1}{p{\gnumericColE}|}%
	{\gnumericPB{\centering}\gnumbox{\textbf{0}}}
	&\multicolumn{1}{p{\gnumericColF}|}%
	{\gnumericPB{\centering}\gnumbox{\textbf{1}}}
	&\multicolumn{1}{p{\gnumericColG}|}%
	{\gnumericPB{\centering}\gnumbox{\textbf{1}}}
	&\multicolumn{1}{p{\gnumericColH}|}%
	{\gnumericPB{\centering}\gnumbox{\textbf{1}}}
\\
\hhline{|--------|}
	 \multicolumn{1}{|p{\gnumericColA}|}%
	{\gnumericPB{\centering}\gnumbox{0}}
	&\multicolumn{1}{p{\gnumericColB}|}%
	{\gnumericPB{\centering}\gnumbox{1}}
	&\multicolumn{1}{p{\gnumericColC}|}%
	{\gnumericPB{\centering}\gnumbox{1}}
	&\multicolumn{1}{p{\gnumericColD}|}%
	{\gnumericPB{\centering}\gnumbox{1}}
	&\multicolumn{1}{p{\gnumericColE}|}%
	{\gnumericPB{\centering}\gnumbox{\textbf{0}}}
	&\multicolumn{1}{p{\gnumericColF}|}%
	{\gnumericPB{\centering}\gnumbox{\textbf{1}}}
	&\multicolumn{1}{p{\gnumericColG}|}%
	{\gnumericPB{\centering}\gnumbox{\textbf{1}}}
	&\multicolumn{1}{p{\gnumericColH}|}%
	{\gnumericPB{\centering}\gnumbox{\textbf{0}}}
\\
\hhline{|--------|}
	 \multicolumn{1}{|p{\gnumericColA}|}%
	{\gnumericPB{\centering}\gnumbox{0}}
	&\multicolumn{1}{p{\gnumericColB}|}%
	{\gnumericPB{\centering}\gnumbox{1}}
	&\multicolumn{1}{p{\gnumericColC}|}%
	{\gnumericPB{\centering}\gnumbox{1}}
	&\multicolumn{1}{p{\gnumericColD}|}%
	{\gnumericPB{\centering}\gnumbox{0}}
	&\multicolumn{1}{p{\gnumericColE}|}%
	{\gnumericPB{\centering}\gnumbox{\textbf{0}}}
	&\multicolumn{1}{p{\gnumericColF}|}%
	{\gnumericPB{\centering}\gnumbox{\textbf{1}}}
	&\multicolumn{1}{p{\gnumericColG}|}%
	{\gnumericPB{\centering}\gnumbox{\textbf{0}}}
	&\multicolumn{1}{p{\gnumericColH}|}%
	{\gnumericPB{\centering}\gnumbox{\textbf{1}}}
\\
\hhline{|--------|}
	 \multicolumn{1}{|p{\gnumericColA}|}%
	{\gnumericPB{\centering}\gnumbox{0}}
	&\multicolumn{1}{p{\gnumericColB}|}%
	{\gnumericPB{\centering}\gnumbox{1}}
	&\multicolumn{1}{p{\gnumericColC}|}%
	{\gnumericPB{\centering}\gnumbox{0}}
	&\multicolumn{1}{p{\gnumericColD}|}%
	{\gnumericPB{\centering}\gnumbox{1}}
	&\multicolumn{1}{p{\gnumericColE}|}%
	{\gnumericPB{\centering}\gnumbox{\textbf{0}}}
	&\multicolumn{1}{p{\gnumericColF}|}%
	{\gnumericPB{\centering}\gnumbox{\textbf{1}}}
	&\multicolumn{1}{p{\gnumericColG}|}%
	{\gnumericPB{\centering}\gnumbox{\textbf{0}}}
	&\multicolumn{1}{p{\gnumericColH}|}%
	{\gnumericPB{\centering}\gnumbox{\textbf{0}}}
\\
\hhline{|--------|}
	 \multicolumn{1}{|p{\gnumericColA}|}%
	{\gnumericPB{\centering}\gnumbox{0}}
	&\multicolumn{1}{p{\gnumericColB}|}%
	{\gnumericPB{\centering}\gnumbox{1}}
	&\multicolumn{1}{p{\gnumericColC}|}%
	{\gnumericPB{\centering}\gnumbox{0}}
	&\multicolumn{1}{p{\gnumericColD}|}%
	{\gnumericPB{\centering}\gnumbox{0}}
	&\multicolumn{1}{p{\gnumericColE}|}%
	{\gnumericPB{\centering}\gnumbox{\textbf{0}}}
	&\multicolumn{1}{p{\gnumericColF}|}%
	{\gnumericPB{\centering}\gnumbox{\textbf{0}}}
	&\multicolumn{1}{p{\gnumericColG}|}%
	{\gnumericPB{\centering}\gnumbox{\textbf{1}}}
	&\multicolumn{1}{p{\gnumericColH}|}%
	{\gnumericPB{\centering}\gnumbox{\textbf{1}}}
\\
\hhline{|--------|}
	 \multicolumn{1}{|p{\gnumericColA}|}%
	{\gnumericPB{\centering}\gnumbox{0}}
	&\multicolumn{1}{p{\gnumericColB}|}%
	{\gnumericPB{\centering}\gnumbox{0}}
	&\multicolumn{1}{p{\gnumericColC}|}%
	{\gnumericPB{\centering}\gnumbox{1}}
	&\multicolumn{1}{p{\gnumericColD}|}%
	{\gnumericPB{\centering}\gnumbox{1}}
	&\multicolumn{1}{p{\gnumericColE}|}%
	{\gnumericPB{\centering}\gnumbox{\textbf{0}}}
	&\multicolumn{1}{p{\gnumericColF}|}%
	{\gnumericPB{\centering}\gnumbox{\textbf{0}}}
	&\multicolumn{1}{p{\gnumericColG}|}%
	{\gnumericPB{\centering}\gnumbox{\textbf{1}}}
	&\multicolumn{1}{p{\gnumericColH}|}%
	{\gnumericPB{\centering}\gnumbox{\textbf{0}}}
\\
\hhline{|--------|}
	 \multicolumn{1}{|p{\gnumericColA}|}%
	{\gnumericPB{\centering}\gnumbox{0}}
	&\multicolumn{1}{p{\gnumericColB}|}%
	{\gnumericPB{\centering}\gnumbox{0}}
	&\multicolumn{1}{p{\gnumericColC}|}%
	{\gnumericPB{\centering}\gnumbox{1}}
	&\multicolumn{1}{p{\gnumericColD}|}%
	{\gnumericPB{\centering}\gnumbox{0}}
	&\multicolumn{1}{p{\gnumericColE}|}%
	{\gnumericPB{\centering}\gnumbox{\textbf{0}}}
	&\multicolumn{1}{p{\gnumericColF}|}%
	{\gnumericPB{\centering}\gnumbox{\textbf{0}}}
	&\multicolumn{1}{p{\gnumericColG}|}%
	{\gnumericPB{\centering}\gnumbox{\textbf{0}}}
	&\multicolumn{1}{p{\gnumericColH}|}%
	{\gnumericPB{\centering}\gnumbox{\textbf{1}}}
\\
\hhline{|--------|}
	 \multicolumn{1}{|p{\gnumericColA}|}%
	{\gnumericPB{\centering}\gnumbox{0}}
	&\multicolumn{1}{p{\gnumericColB}|}%
	{\gnumericPB{\centering}\gnumbox{0}}
	&\multicolumn{1}{p{\gnumericColC}|}%
	{\gnumericPB{\centering}\gnumbox{0}}
	&\multicolumn{1}{p{\gnumericColD}|}%
	{\gnumericPB{\centering}\gnumbox{1}}
	&\multicolumn{1}{p{\gnumericColE}|}%
	{\gnumericPB{\centering}\gnumbox{\textbf{0}}}
	&\multicolumn{1}{p{\gnumericColF}|}%
	{\gnumericPB{\centering}\gnumbox{\textbf{0}}}
	&\multicolumn{1}{p{\gnumericColG}|}%
	{\gnumericPB{\centering}\gnumbox{\textbf{0}}}
	&\multicolumn{1}{p{\gnumericColH}|}%
	{\gnumericPB{\centering}\gnumbox{\textbf{0}}}
\\
\hhline{|-|-|-|-|-|-|-|-|}
%\end{longtable}
\end{tabular}




\ifthenelse{\isundefined{\languageshorthands}}{}{\languageshorthands{\languagename}}
\gnumericTableEnd

	\caption{}
	\label{table:counter_decoder}
\end{table}
\section{Karnaugh Map}
Using Boolean logic, output $A$  in Table \ref{table:counter_decoder} can be expressed in terms of the inputs $W,X,Y,Z$ as
\begin{equation}
\label{eq:A}
A = W^{\prime}X^{\prime}Y^{\prime}Z^{\prime}+W^{\prime}XY^{\prime}Z^{\prime}\\
+W^{\prime}X^{\prime}YZ^{\prime}+W^{\prime}XYZ^{\prime}+W^{\prime}X^{\prime}Y^{\prime}Z
\end{equation}
\begin{enumerate}
\item K-Map for $A$: 
The expression in \eqref{eq:A}  can be minimized using the K-map in Fig. \ref{fig:kmap_A}.
In Fig. \ref{fig:kmap_A},  the {\em implicants} in boxes 0,2,4,6 result in $W^{\prime}Z^{\prime}$.  The implicants in
boxes 0,8 result in $W^{\prime}X^{\prime}Y^{\prime}$.  Thus, after minimization using Fig. \ref{eq:kmap_A},  \eqref{eq:A} can be expressed as
\begin{equation}
	\label{eq:kmap_A}
	A = W^{\prime}Z^{\prime}+W^{\prime}X^{\prime}Y^{\prime}
\end{equation}
%
Using the fact that
\begin{align}
	\label{eq:boolean}
	\begin{split}
		X+X^{\prime} &= 1
		\\
		XX^{\prime} &= 0,
	\end{split}
\end{align}
%
derive \eqref{eq:kmap_A} from \eqref{eq:A} algebraically.\\
\textbf{Solution: } 
\begin{equation}
	A = W^{\prime}X^{\prime}Y^{\prime}Z^{\prime}+W^{\prime}XY^{\prime}Z^{\prime}+W^{\prime}X^{\prime}YZ^{\prime}+W^{\prime}XYZ^{\prime}+W^{\prime}X^{\prime}Y^{\prime}Z 
\end{equation}
\begin{equation}
	A = W^{\prime}X^{\prime}Y^{\prime}+W^{\prime}X^{\prime}Z^{\prime}
	+W^{\prime}Y^{\prime}Z^{\prime}+W^{\prime}XYZ^{\prime}\\
\end{equation}
\begin{equation}
	A = W^{\prime}X^{\prime}Y^{\prime}+W^{\prime}Z^{\prime}(X^{\prime}+Y^{\prime}+XY)\\
\end{equation}
\begin{equation}
	A = W^{\prime}X^{\prime}Y^{\prime}+W^{\prime}Z^{\prime}((XY)^{\prime}+XY)\\
\end{equation}
\begin{equation}
		A = W^{\prime}Z^{\prime}+W^{\prime}X^{\prime}Y^{\prime}	
\end{equation}

%
%
%
\begin{figure}[!h]
	\resizebox {\columnwidth} {!} {
		\begin{karnaugh-map}[4][4][1][][]
    \maxterms{1,3,5,7,9,10,11,12,13,14,15}
    \minterms{0,2,4,6,8}
    \implicantedge{0}{4}{2}{6}
    \implicantedge{0}{0}{8}{8}
    % note: posistion for start of \draw is (0, Y) where Y is
    % the Y size(number of cells high) in this case Y=2
    \draw[color=black, ultra thin] (0, 4) --
    node [pos=0.7, above right, anchor=south west] {$XW$} % Y label
    node [pos=0.7, below left, anchor=north east] {$ZY$} % X label
    ++(135:1);
        
    \end{karnaugh-map}

	}
	\caption{K-map for $A$.}
	\label{fig:kmap_A}
\end{figure}
%
\item K-Map for $B$:
From Table \ref{table:counter_decoder}, using boolean logic,
\begin{equation}
	\label{eq:B}
	B = WX^{\prime}Y^{\prime}Z^{\prime} + W^{\prime}XY^{\prime}Z^{\prime}
	+WX^{\prime}YZ^{\prime}
	+W^{\prime}XYZ^{\prime}
\end{equation}
%
\begin{figure}[!h]
	\resizebox {\columnwidth} {!} {
		\begin{karnaugh-map}[4][4][1][][]
    \maxterms{0,3,4,7,8,9,10,11,12,13,14,15}
    \minterms{1,2,5,6}
     \implicant{1}{5}    
     \implicant{2}{6}         
    % note: posistion for start of \draw is (0, Y) where Y is
    % the Y size(number of cells high) in this case Y=2
    \draw[color=black, ultra thin] (0, 4) --
    node [pos=0.7, above right, anchor=south west] {$XW$} % Y label
    node [pos=0.7, below left, anchor=north east] {$ZY$} % X label
    ++(135:1);
        
    \end{karnaugh-map}

	}
	\caption{K-map for $B$.}
	\label{fig:kmap_B}
\end{figure}
%
Show that \eqref{eq:B} can be reduced to
\begin{equation}
	\label{eq:kmap_B}
	B = WX^{\prime}Z^{\prime} + W^{\prime}XZ^{\prime}
\end{equation}
using Fig. \ref{fig:kmap_B}.
\item Derive \eqref{eq:kmap_B} from \eqref{eq:B} algebraically using \eqref{eq:boolean}.
\textbf{Solution:}
\begin{equation}
	\label{eq:B}
	B = WX^{\prime}Y^{\prime}Z^{\prime} + W^{\prime}XY^{\prime}Z^{\prime}
	+WX^{\prime}YZ^{\prime}
	+W^{\prime}XYZ^{\prime}
\end{equation}
\begin{equation}
	B = WX^{\prime}Z^{\prime} + W^{\prime}XZ^{\prime}
\end{equation}
%
\item {K-Map for $C$: }
From Table \ref{table:counter_decoder}, using boolean logic,
\begin{equation}
	\label{eq:C}
	C = WXY^{\prime}Z^{\prime} + W^{\prime}X^{\prime}YZ^{\prime}
	+WX^{\prime}YZ^{\prime}
	+W^{\prime}XYZ^{\prime}
\end{equation}
%
%
\begin{figure}[!h]
	\resizebox {\columnwidth} {!} {
		\begin{karnaugh-map}[4][4][1][][]
    \maxterms{0,1,2,7,8,9,10,11,12,13,14,15}
    \minterms{3,4,5,6}
     \implicant{4}{5}    
     \implicant{3}{3}
    \implicantedge{4}{4}{6}{6}     
%     \implicant{2}{6}         
    % note: posistion for start of \draw is (0, Y) where Y is
    % the Y size(number of cells high) in this case Y=2
    \draw[color=black, ultra thin] (0, 4) --
    node [pos=0.7, above right, anchor=south west] {$XW$} % Y label
    node [pos=0.7, below left, anchor=north east] {$ZY$} % X label
    ++(135:1);
        
    \end{karnaugh-map}

	}
	\caption{K-map for $C$.}
	\label{fig:kmap_C}
\end{figure}
%
Show that \eqref{eq:C} can be reduced to
\begin{equation}
	\label{eq:kmap_C}
	C = WXY^{\prime}Z^{\prime}  +  X^{\prime}YZ^{\prime} + W^{\prime}YZ^{\prime}
\end{equation}
using Fig. \ref{fig:kmap_C}.
%
\item 
Derive \eqref{eq:kmap_C} from \eqref{eq:C} algebraically using \eqref{eq:boolean}.
\textbf{Solution:}
\begin{equation}
	C = WXY^{\prime}Z^{\prime} + W^{\prime}X^{\prime}YZ^{\prime}
	+WX^{\prime}YZ^{\prime}
	+W^{\prime}XYZ^{\prime}
\end{equation}
\begin{equation}
	C = WXY^{\prime}Z^{\prime}  +  X^{\prime}YZ^{\prime} + W^{\prime}YZ^{\prime}
\end{equation}
%
\item {K-Map for $D$: }
From Table \ref{table:counter_decoder}, using boolean logic,
\begin{equation}
	\label{eq:D}
	D = WXYZ^{\prime} + W^{\prime}X^{\prime}Y^{\prime}Z
\end{equation}
%
\begin{figure}[!h]
	\resizebox {\columnwidth} {!} {
		\begin{karnaugh-map}[4][4][1][][]
    \maxterms{0,1,2,3,4,5,6,9,10,11,12,13,14,15}
    \minterms{7,8}
     \implicant{7}{7}    
     \implicant{8}{8}
    % note: posistion for start of \draw is (0, Y) where Y is
    % the Y size(number of cells high) in this case Y=2
    \draw[color=black, ultra thin] (0, 4) --
    node [pos=0.7, above right, anchor=south west] {$XW$} % Y label
    node [pos=0.7, below left, anchor=north east] {$ZY$} % X label
    ++(135:1);
        
    \end{karnaugh-map}

	}
	\caption{K-map for $D$.}
	\label{fig:kmap_D}
\end{figure}
%
\item 
Minimize \eqref{eq:D} using Fig. \ref{fig:kmap_D}.
%
\item Download the code in
\begin{lstlisting}
	wget https://raw.githubusercontent.com/gadepall/arduino/master/7447/codes/inc_dec/inc_dec.ino
\end{lstlisting}
%
and modify it using the K-Map equations for A,B,C and D. Execute and verify.
\item {Display Decoder:}
Table \ref{table:disp_dec} is the truth table for the display decoder.  
Use K-maps to obtain the minimized expressions for $a,b,c,d,e,f,g$ in terms of $A,B,C,D$ without don't care conditions.
%

	\begin{table}[h!]
	\begin{center}
		\begin{tabular}{ |c|c|c|c|c|c|c|c|c|c|c|c| } 
			\hline
			D & C & B & A & \textbf{a} & \textbf{b} & \textbf{c} & \textbf{d} & \textbf{e} & \textbf{f} & \textbf{g} & Decimal \\ 
			\hline
			0 & 0 & 0 & 0 & 0 & 0 & 0 & 0 & 0 & 0 & 1 & 0 \\ 
			\hline
			0 & 0 & 0 & 1 & 1 & 0 & 0 & 1 & 1 & 1 & 1 & 1 \\ 
			\hline
			0 & 0 & 1 & 0 & 0 & 0 & 1 & 0 & 0 & 1 & 0 & 2 \\ 
			\hline
			0 & 0 & 1 & 1 & 0 & 0 & 0 & 0 & 1 & 1 & 0 & 3 \\
			\hline
			0 & 1 & 0 & 0 & 0 & 1 & 0 & 0 & 1 & 0 & 0 & 4 \\
			\hline
			0 & 1 & 0 & 1 & 1 & 0 & 0 & 1 & 1 & 0 & 0 & 5 \\
			\hline
			0 & 1 & 1 & 0 & 0 & 1 & 0 & 0 & 1 & 0 & 0 & 6 \\ 
			\hline
			0 & 1 & 1 & 1 & 0 & 0 & 0 & 1 & 1 & 1 & 1 & 7 \\
			\hline
			1 & 0 & 0 & 0 & 0 & 0 & 0 & 0 & 0 & 0 & 0 & 8 \\
			\hline
			1 & 0 & 0 & 1 & 0 & 0 & 0 & 0 & 1 & 0 & 0 & 9 \\ 
			\hline
		\end{tabular}
		\caption{Truth table for display decoder.}
		\label{table:disp_dec}
	\end{center}
\end{table}
\end{enumerate}
\textbf{Solution:}\\
Without DON'T CARE: \\
from Fig. \ref{fig:kmap_a}
\begin{equation}
	a = D^{\prime}C^{\prime}B^{\prime}A+D^{\prime}CB^{\prime}A^{\prime}
\end{equation}
from Fig. \ref{fig:kmap_b}
\begin{equation}
	b = D^{\prime}CB^{\prime}A+D^{\prime}CBA^{\prime}
\end{equation}
from Fig.  \ref{fig:kmap_c}
\begin{equation}
	c = D^{\prime}C^{\prime}BA^{\prime}
\end{equation}
from Fig. \ref{fig:kmap_d}
\begin{equation}
	d = D^{\prime}CB^{\prime}A^{\prime}+D^{\prime}CBA+C^{\prime}B^{\prime}A
\end{equation}
from Fig. \ref{fig:kmap_e}
\begin{equation}
	e = D^{\prime}A+C^{\prime}B^{\prime}A+D^{\prime}CB^{\prime}
\end{equation}
from Fig. \ref{fig:kmap_f}
\begin{equation}
	f = D^{\prime}BA+D^{\prime}C^{\prime}A+D^{\prime}C^{\prime}B
\end{equation}
from Fig. \ref{fig:kmap_g}
\begin{equation}
	g = D^{\prime}C^{\prime}B^{\prime}+D^{\prime}CBA
\end{equation}
\begin{figure}[!h]
	\resizebox {\columnwidth} {!} {
		\begin{karnaugh-map}[4][4][1][][]
   \maxterms{0,2,3,5,6,7,8,9}
    \minterms{1,4}
    \implicant{4}{12}
	\indeterminants{10,11,12,13,14,15}        
    \draw[color=black, ultra thin] (0, 4) --
    node [pos=0.7, above right, anchor=south west] {$BA$} % Y label
    node [pos=0.7, below left, anchor=north east] {$DC$} % X label
    ++(135:1);
        
    \end{karnaugh-map}

	}
	\caption{K-map for $a$.}
	\label{fig:kmap_a}
\end{figure}
\begin{figure}[!h]
	\resizebox {\columnwidth} {!} {
		\begin{karnaugh-map}[4][4][1][][]
    \maxterms{0,1,2,5,6,9,10,11,12,13,14,15}
    \minterms{3,4,7,8}
    \implicant{3}{7}
    \implicant{4}{4}
    \implicant{8}{8}
    % note: posistion for start of \draw is (0, Y) where Y is
    % the Y size(number of cells high) in this case Y=2
    \draw[color=black, ultra thin] (0, 4) --
    node [pos=0.7, above right, anchor=south west] {$XW$} % Y label
    node [pos=0.7, below left, anchor=north east] {$ZY$} % X label
    ++(135:1);
        
    \end{karnaugh-map}

	}
	\caption{K-map for $b$.}
	\label{fig:kmap_b}
\end{figure}
\begin{figure}[!h]
	\resizebox {\columnwidth} {!} {
		\begin{karnaugh-map}[4][4][1][][]
    \maxterms{0,1,2,3,4,9,10,11,12,13,14,15}
    \minterms{5,6,7,8}
    \implicant{5}{7}
    \implicant{7}{6}
    \implicant{8}{8}
    % note: posistion for start of \draw is (0, Y) where Y is
    % the Y size(number of cells high) in this case Y=2
    \draw[color=black, ultra thin] (0, 4) --
    node [pos=0.7, above right, anchor=south west] {$XW$} % Y label
    node [pos=0.7, below left, anchor=north east] {$ZY$} % X label
    ++(135:1);
        
    \end{karnaugh-map}

	}
	\caption{K-map for $c$.}
	\label{fig:kmap_c}
\end{figure}
\begin{figure}[!h]
	\resizebox {\columnwidth} {!} {
		\begin{karnaugh-map}[4][4][1][][]
    \maxterms{0,2,3,5,6,8,10,11,12,13,14,15}
    \minterms{1,4,7,9}
    \implicant{7}{7}
    \implicant{4}{4}
    \implicantedge{1}{1}{9}{9}
    % note: posistion for start of \draw is (0, Y) where Y is
    % the Y size(number of cells high) in this case Y=2
    \draw[color=black, ultra thin] (0, 4) --
    node [pos=0.7, above right, anchor=south west] {$BA$} % Y label
    node [pos=0.7, below left, anchor=north east] {$DC$} % X label
    ++(135:1);
        
    \end{karnaugh-map}

	}
	\caption{K-map for $d$.}
	\label{fig:kmap_d}
\end{figure}
\begin{figure}[!h]
	\resizebox {\columnwidth} {!} {
		\begin{karnaugh-map}[4][4][1][][]
    \maxterms{0,2,6,8,10,11,12,13,14,15}
    \minterms{1,3,4,5,7,9}
    \implicantedge{1}{1}{9}{9}
    \implicant{4}{5}
    \implicant{1}{7}
    % note: posistion for start of \draw is (0, Y) where Y is
    % the Y size(number of cells high) in this case Y=2
    \draw[color=black, ultra thin] (0, 4) --
    node [pos=0.7, above right, anchor=south west] {$BA$} % Y label
    node [pos=0.7, below left, anchor=north east] {$DC$} % X label
    ++(135:1);
        
    \end{karnaugh-map}

	}
	\caption{K-map for $e$.}
	\label{fig:kmap_e}
\end{figure}
\begin{figure}[!h]
	\resizebox {\columnwidth} {!} {
		\begin{karnaugh-map}[4][4][1][][]
   \maxterms{0,4,5,6,8,9}
    \minterms{1,2,3,7}
    \implicantedge{3}{2}{11}{10}
    \implicant{3}{11}
    \implicant{1}{3}
	\indeterminants{10,11,12,13,14,15}        
    \draw[color=black, ultra thin] (0, 4) --
    node [pos=0.7, above right, anchor=south west] {$BA$} % Y label
    node [pos=0.7, below left, anchor=north east] {$DC$} % X label
    ++(135:1);
        
    \end{karnaugh-map}

	}
	\caption{K-map for $f$.}
	\label{fig:kmap_f}
\end{figure}
\begin{figure}[!h]
	\resizebox {\columnwidth} {!} {
		\begin{karnaugh-map}[4][4][1][][]
   \maxterms{2,3,4,5,6,8,9}
    \minterms{0,1,7}
    \implicant{0}{1}
    \implicant{7}{15}
	\indeterminants{10,11,12,13,14,15}        
    \draw[color=black, ultra thin] (0, 4) --
    node [pos=0.7, above right, anchor=south west] {$BA$} % Y label
    node [pos=0.7, below left, anchor=north east] {$DC$} % X label
    ++(135:1);
        
    \end{karnaugh-map}

	}
	\caption{K-map for $g$.}
	\label{fig:kmap_g}
\end{figure}
\end{document}


